\section{尾部风险与压力测试}
\label{ch:lower-portfolio-tail}

\begin{itemize}
  % Generated from curriculum/manifest.json; do not edit by hand.
\item 直接先修：下册《组合优化：从观点后验到稳健实施》。

\end{itemize}

分位数只给一个阈值，期望短缺还要把超过阈值的概率质量计入平均。离散样本的边界质量最容易算错，因此先做一个非整数尾部个数的例子。


\MFQTransition{从分位数到尾部平均}

对损失 $L=-R_p$，左分位数定义为
\[
 \operatorname{VaR}_\alpha(L)=\inf\{\ell:\mathbb P(L\le\ell)\ge\alpha\},
\]
积分分位数 ES 为
\[
 \operatorname{ES}_\alpha(L)=\frac1{1-\alpha}
 \int_\alpha^1\operatorname{VaR}_u(L)\,\mathrm du.
\]
离散样本在分位点可能有概率质量，因此软件默认的线性插值、高位次序统计量和左分位数会得到不同答案。本节沿用上述经验左分位数定义，下面直接按概率质量计算。

\MFQTransition{离散 ES 的边界质量}

四个等概率损失为 $(0,1,2,10)$，$\alpha=0.6$。左分位数为 2，因为损失不超过 1 的概率是 0.5，不超过 2 的概率是 0.75。最坏 40\% 包含损失 10 的全部 25\% 质量，以及损失 2 的 15\% 质量，故
\[
 \operatorname{ES}_{0.6}=\frac{0.25(10)+0.15(2)}{0.4}=7.
\]
直接取所有 $L\geq2$ 的条件平均得到 6，是最坏 50\% 的平均，不是这里的 ES。代入辅助阈值公式 $z+\mathbb E(L-z)_+/(1-\alpha)$，在 $z=2$ 得 $2+(0.25\cdot8)/0.4=7$，两条计算一致。

% Generated by tools/build_figures.py; edit figures/figure-specs.json instead.
\MFQFigure{tex/figures/generated/lower-pf-tail.pdf}{VaR 给出分位点，ES 按尾部概率平均损失；分位点上的概率质量有时只计入一部分。尾部观察稀少时，估计的不确定性较大。}{lower-pf-tail}


VaR 一般不满足次可加性；ES 在常见可积条件下是相干风险度量，但更依赖尾部样本。关于市场、信用与操作风险中的标准定义、估计与回测框架，可参见 McNeil、Frey 与 Embrechts\cite{mcneil2015quantitative}。有效尾部观察数为
\[
 n_{\mathrm{tail}}=n(1-\alpha).
\]
例如 250 期观测在 99\% 水平只对应约 2.5 个尾部观测。没有一个普遍适用的“达到 20 或 30 就可信”定理；可信度还取决于尾部形状、依赖和持仓稳定性。读者应比较置信水平和样本窗口变化，并解释少量极端观测对结果的影响。

\MFQTransition{历史尾部、参数尾部和压力测试回答不同问题}

历史法问“过去样本中尾部如何”；参数法问“若分布模型成立，尾部如何”；压力测试问“指定联合冲击下损失多少”。三者不能互相替代。轻尾正态可能漏掉跳跃、波动聚集和相关性上升；历史样本可能没有经历当前持仓结构；压力情景则没有自动的发生概率。

联合情景描述哪些价格和流动性变化同时发生，并将它们作用于同一时点的持仓。历史重演保留曾发生的联合冲击，假设情景改变利率、波动率、价差或相关性，反向压力则从资本或流动性阈值反推触发条件。

\MFQTransition{非线性头寸不能只乘 Delta}

设每个风险因子的线性暴露为 $d_i$、二阶暴露为 $g_i$，冲击为 $x_i$，二阶损益近似
\[
 \Delta V(x)=d^\top x+\frac12\sum_i g_ix_i^2,
 \qquad L(x)=-\Delta V(x).
\]
取冲击方向 $x=(-0.10,0.04)$、$d=(60,40)$、$g=(80,-20)$。线性损失为 $4.4$，加入二阶项后为 $4.016$。这说明非线性重估不一定单调放大损失；Gamma 符号、交叉项、波动率曲面和路径依赖都会改变结果。

完整重估在新的风险因子取值下重新计算价值，适合检验二阶近似在大冲击下的误差。期权需要现货、波动率、时间和利率共同变化；债券需要曲线重构；信用与流动性头寸还需要违约、价差和成交假设。含期权时还会出现 $\mathrm{Vega}\,\Delta\sigma$ 与交叉 Gamma 项。大冲击、障碍附近或临近到期时，Taylor 近似可能失效，应直接重算估值；各风险因子的冲击还需能共同构成一个有意义的情景。

\MFQTransition{反向压力寻找的是第一处阈值穿越}

沿方向 $x$ 放大尺度 $s\ge0$，反向压力问题为
\[
 s^*=\inf\{s:L(sx)\ge C\},
\]
其中 $C$ 是资本或损失阈值。阈值取 $C=10$，需要求解二次不等式。

代入本章给定暴露与冲击，二阶损失为 $L(s)=4.4s-0.384s^2$，阈值 10 的两根为 $3.125$ 与 $25/3$。从零增加压力时先碰到 $3.125$；只检查很大的 $s$，反而可能看到损失重新低于阈值。二阶近似在大冲击下本就可能不可靠，因此求得根后还要判断近似适用范围。若只在某个方向放大，结果仅是该方向的阈值；更广泛的反向压力需要另给允许方向与幅度。Notebook 画出这条曲线，并允许改变暴露与阈值。


\MFQTransition{从样本误差到模型误差}

样本较短时，尾部平均主要由少数极端观测决定。可用块 bootstrap 保留一部分时间依赖，比较重抽样得到的风险估计；区间本身仍受块长度和未观察尾部的限制。极值理论用高阈值以上的超额拟合尾部分布，但阈值太高使数据稀少，太低又使尾部近似不合适，形成偏差与方差的权衡。

VaR 的样本外评价可记录实际损失超过预测分位点的次数：若条件分位点模型正确且分布连续，每期超损的条件概率应为 $1-\alpha$。只比较总次数还不够，连续聚集的超损可能提示条件波动未被解释。ES 还涉及这些极端损失的大小，因此不能仅凭超损次数评价尾部平均。

\section{本路线的综合项目}

把协方差估计、优化约束和尾部压力放在同一个组合中，比较历史、假设和反向压力情景。报告应同时给出风险贡献、成本、约束的影子价格和失败阈值，并说明哪些结果只依赖当前样本，哪些可以迁移到下一段数据。

\section{分层练习}
\begin{enumerate}[leftmargin=*]
  \item \textbf{口述概念：}为什么协方差半正定仍不足以保证估计可靠？
  \item \textbf{笔试推导：}证明方差成分贡献之和等于组合方差。
  \item \textbf{笔试推导：}对角协方差下推导风险平价权重与逆波动率成正比。
  \item \textbf{数值编程：}改变收缩强度，报告最小特征值、条件数和风险平价权重。
  \item \textbf{数值编程：}把 IID bootstrap 改为长度为 4 的移动区块 bootstrap。
  \item \textbf{研究审计：}找出成对删除导致非半正定矩阵的最小例子。
  \item \textbf{研究判断：}因子模型何时比样本协方差更稳定，却更可能遗漏风险？
  \item \textbf{开放任务：}为真实收益数据写出资产池、时间索引、公司行动与缺失值协议。
\end{enumerate}

\section{分层练习}
\begin{enumerate}[leftmargin=*]
  \item \textbf{口述概念：}Black--Litterman 中 $P,q,\Omega,\tau$ 分别编码什么？
  \item \textbf{笔试推导：}从高斯精度相加推导后验均值。
  \item \textbf{笔试推导：}从 $(L-z)_+$ 推导 CVaR 的辅助变量 LP。
  \item \textbf{数值编程：}用情景枚举核对 LP 的权重、阈值和目标值。
  \item \textbf{数值编程：}把最大权重从 0.8 改为 0.6，重算 CVaR 和活跃约束。
  \item \textbf{错误研究包审计：}找出“先优化再裁剪权重”破坏的约束与目标。
  \item \textbf{研究判断：}比较均值不确定性惩罚、协方差 ridge 和交易成本的经济含义。
  \item \textbf{开放任务：}加入部分成交与整数头寸，写出恢复协议和独立现金台账。
\end{enumerate}

\MFQTransition{尾部、压力测试与模型风险题组}
\begin{enumerate}[leftmargin=*]
  \item \textbf{口述概念：}VaR、ES、历史压力和反向压力分别回答什么问题？
  \item \textbf{笔试推导：}解释积分分位数 ES 如何处理分位点概率质量。
  \item \textbf{笔试推导：}推导二阶重估下反向压力的标量方程。
  \item \textbf{数值编程：}对 500 个样本比较 95\% 与 99\% 下的尾部观察数，并解释估计敏感性。
  \item \textbf{数值编程：}加入交叉 Gamma 或期权全重估，比较线性与非线性损失。
  \item \textbf{错误研究材料审计：}审查只报告四个损失样本 99\% ES 的报告。
  \item \textbf{研究判断：}为什么反向压力阈值不等于危机发生概率？
\end{enumerate}

四个样本估计 99\% ES 时，最坏 1\% 概率质量完全落在最大观察上。请说明这一计算在经验分布下的含义，以及为什么它没有提供未观察尾部的信息。
\subsection*{基础衔接：独立计算}
对同样四个损失，求 $\alpha=0.75$ 的 VaR 与 ES。

